
\section{Architecture globale}

La solution est organisée en quatre ensembles : exposition, couche applicative, services d'intelligence artificielle et stockages. Les pipelines documentaires alimentent Qdrant indépendamment du traitement conversationnel en ligne.

\begin{figure}[H]
\centering
\begin{tikzpicture}[
  node distance=8mm and 12mm,
  box/.style={draw, rounded corners, align=center, minimum width=3.1cm, minimum height=1.1cm},
  flow/.style={-Latex, thick}
]
\node[box] (user) {Collaborateur\\Navigateur};
\node[box, right=of user] (nginx) {Nginx\\TLS / reverse proxy};
\node[box, right=of nginx] (ui) {Streamlit\\SSO et interface};
\node[box, right=of ui] (backend) {FastAPI\\Orchestration RAG};
\node[box, below left=of backend] (embed) {Embeddings\\multilingual-e5-large};
\node[box, below=of backend] (qdrant) {Qdrant\\Vecteurs et payloads};
\node[box, below right=of backend] (llm) {vLLM\\Mistral Small 3.1 24B};
\node[box, right=of backend] (pg) {PostgreSQL\\Interactions};
\draw[flow] (user) -- (nginx);
\draw[flow] (nginx) -- (ui);
\draw[flow] (ui) -- (backend);
\draw[flow] (backend) -- (embed);
\draw[flow] (backend) -- (qdrant);
\draw[flow] (backend) -- (llm);
\draw[flow] (backend) -- (pg);
\end{tikzpicture}
\caption{Architecture logique de l'Assistant RAG DNSI}
\label{fig:architecture_globale_ch5}
\end{figure}

\subsection{Services de production}

\begin{table}[H]
\centering
\caption{Services de la stack active}
\label{tab:services_ch5}
\begin{tabular}{p{2.0cm}p{2.4cm}p{1.8cm}p{4.4cm}p{3.1cm}}
\toprule
\textbf{Service} & \textbf{Technologie} & \textbf{Port} & \textbf{Rôle} & \textbf{Point d'attention} \\
\midrule
Nginx & nginx:1.27-alpine & 80/443 & Terminaison TLS et reverse proxy vers l’interface. & Certificats montés en lecture seule. \\
Interface & Streamlit & 8502 & SSO, saisie de la question, historique, sources, visualisation PDF et feedback. & Dépend du backend sain. \\
Backend & FastAPI & 8000 interne & Orchestration RAG, routes, ACL, persistance et gestion de charge. & Non exposé directement par le Compose de production. \\
Embeddings & FastAPI + SentenceTransformers & 8001 & Encodage query/passage avec multilingual-e5-large et normalisation L2. & Modèle monté depuis /models. \\
LLM & vLLM OpenAI-compatible & 8000 interne ; 8002 loopback & Inférence Mistral Small 3.1 24B. & GPU requis ; clé API ; healthcheck. \\
Qdrant & qdrant/qdrant & 6333/6334 loopback & Collections vectorielles, payloads, recherche et filtrage. & Persistance dans qdrant\_data. \\
PostgreSQL & postgres:16-alpine & 5432 interne & Sessions, interactions, sources, métadonnées et feedbacks. & Volume pgdata ; mot de passe obligatoire. \\
\bottomrule
\end{tabular}
\end{table}

\subsection{Topologie réseau}

Tous les services appartiennent au réseau Docker bridge \texttt{chatbot\_net}. Le backend n'est pas publié directement sur l'hôte dans la configuration de production ; l'interface l'atteint par son nom de service. vLLM et Qdrant sont publiés sur l'interface loopback, tandis que Nginx expose les ports web.

\begin{table}[H]
\centering
\caption{Contrats réseau principaux}
\label{tab:ports_ch5}
\begin{tabular}{p{3.5cm}p{3.0cm}p{2.0cm}p{5.0cm}}
\toprule
\textbf{Flux} & \textbf{Protocole} & \textbf{Port} & \textbf{Finalité} \\
\midrule
Navigateur -> Nginx & HTTPS & 443 & Accès utilisateur. \\
Nginx -> Streamlit & HTTP & 8502 & Routage interne de l’interface. \\
Streamlit -> FastAPI & HTTP/JSON & 8000 & Routes de santé, session, chat, historique et feedback. \\
FastAPI -> Embeddings & HTTP/JSON & 8001 & Encodage des requêtes et passages. \\
FastAPI -> Qdrant & HTTP/gRPC client & 6333/6334 & Recherche vectorielle et filtres de payload. \\
FastAPI -> vLLM & HTTP compatible OpenAI & 8000 & Appel /v1/chat/completions. \\
FastAPI -> PostgreSQL & TCP/SQL & 5432 & Persistance applicative. \\
\bottomrule
\end{tabular}
\end{table}

\section{Flux d'une requête RAG}

\begin{enumerate}
\item L’interface récupère l’identité et les groupes disponibles dans les claims SSO, puis construit la requête de chat.
\item Le client HTTP envoie au backend la question, l’historique, l’identifiant de conversation, le périmètre documentaire et les principaux disponibles.
\item FastAPI vérifie la clé d’API, applique la limite globale de requêtes et qualifie la question avec le classifieur.
\item Lorsque la recherche documentaire est nécessaire, le backend demande au service d’embeddings d’encoder la requête avec la convention E5 « query: ».
\item Le backend construit le filtre ACL à partir des principaux normalisés et interroge une ou plusieurs collections Qdrant.
\item Les résultats sont harmonisés, éventuellement complétés par un score lexical, rerankés et réduits aux extraits utiles.
\item Le prompt système impose l’utilisation des extraits autorisés et définit le comportement d’abstention lorsque l’information n’est pas présente.
\item Le client vLLM appelle l’endpoint /v1/chat/completions avec le modèle servi et ajuste le budget de sortie lorsque la fenêtre de contexte l’exige.
\item La réponse, les sources et les métadonnées sont enregistrées dans PostgreSQL avant d’être renvoyées à l’interface.
\item Streamlit affiche la réponse, les sources et le visualiseur documentaire, puis permet d’associer un feedback à l’interaction.
\end{enumerate}

Le backend est le point d'orchestration. L'interface n'appelle directement ni Qdrant ni vLLM dans le chemin de production décrit.
